% ============================================================================
% BU-COMPLIANT BUILD OF THE STANDALONE RESUME
%
% Source of truth: ~/Desktop/resume-2025/resume
% This wrapper re-typesets that resume with the margins BU requires, and is
% compiled by ./build-cv.sh into ../graphics/cv.pdf, which endmatter/cv.tex
% then includes as the Curriculum Vitae of the dissertation.
%
% Why not just \includepdf the resume's own main_resume.pdf?
% Its geometry is left/right 0.75in and top/bottom 0.5in, and it sets
% \pagestyle{empty}. BU requires top 1.5in, left 1.5in, right 1in, bottom 1in
% and a page number on every page. Scaling the existing PDF down to fit would
% shrink its 11pt body text to roughly 9.4pt, below the 10pt floor set by
% BU guide 1.1.2 -- so it is re-typeset at full size instead.
%
% Everything else (layout, burgundy headings, FontAwesome icons, section
% order) is inherited unchanged from the resume sources, so editing the resume
% and re-running ./build-cv.sh updates the dissertation.
% ============================================================================

\documentclass{resume}

% Override resume.cls geometry with the BU margins.
% Bottom is 1.1in rather than 1.0in to leave headroom for descenders, and the
% page number sits inside that band, added by pdfpages in endmatter/cv.tex.
\geometry{
	paperwidth  = 8.5in,
	paperheight = 11in,
	left        = 1.5in,
	right       = 1in,
	top         = 1.5in,
	bottom      = 1.1in
}

% Links keep the resume's own accent colour (harvardburgundy, set by
% resume.cls) rather than being blacked out. The contact line now shows LinkedIn
% as a label rather than a spelled-out address, so if the link were not visibly
% a link it would read as the bare word "LinkedIn" with no way to reach it and
% no URL on the printed page.
%
% For a fully monochrome submission, add back:
%     \hypersetup{hidelinks}
% but then also change sections/header.tex to spell out an address, because a
% black unstyled label is unreachable in print.

% resume.cls sets \hyphenpenalty=10000 to stop words breaking across lines.
% A side effect is that long DOI strings cannot break either, and they then run
% past the 1in right margin, which BU guide 1.1.3 forbids outright ("Do not
% exceed the right hand 1-inch margin").
%
% In publications.tex the DOIs are written as \href{<url>}{<same url>}, so the
% visible text is ordinary text rather than \url{...} and xurl alone cannot
% break it. Patch \href so that when the link text is identical to the target
% -- the bare-DOI case -- it is typeset with \nolinkurl, which xurl allows to
% break after any character. Links whose text differs from the target (the lab
% website, the Dolia repository, the Bachelor's thesis) pass through untouched,
% so their formatting and colour survive.
\usepackage{xurl}
\usepackage{etoolbox}
\urlstyle{same} % break like a URL, but keep the surrounding roman font

% The Publications, Lab Skills and Technical Skills sections are built with
% tabularx (see \tableEnv / \tableEnvPlain in resume.cls). A tabularx is a
% single unbreakable box, so when one does not fit in the space left on a page
% it jumps to the next one whole, stranding its section heading above a large
% gap -- at BU margins the Publications table left about 3.8 inches blank.
% ltablex reimplements tabularx on top of longtable so the same tables can
% break across pages; \keepXColumns preserves the auto-sizing X column, which
% ltablex would otherwise downgrade to a plain l column.
\usepackage{ltablex}
\keepXColumns

% Keep a section heading with its content. Without this, a heading that lands
% near the foot of a page is left stranded there while its body starts on the
% next page. Requiring four lines of room forces the heading to move down with
% its content instead. Worth having because the resume changes often and its
% pagination shifts with it.
\usepackage{needspace}
\renewenvironment{rSection}[1]{%
	\needspace{4\baselineskip}%
	\section*{#1}\vspace{-0.4em}\hrule\vspace{0.6em}%
}{\vspace{0.2em}}

\let\BUoriginalHref\href
\renewcommand{\href}[2]{%
	\ifstrequal{#1}{#2}%
		{\BUoriginalHref{#1}{\nolinkurl{#2}}}%
		{\BUoriginalHref{#1}{#2}}%
}

\begin{document}

	% BU guide 1.1.6 names this section the Vita / Curriculum Vitae. Give it a
	% heading so a format reviewer can find it and so the Table of Contents
	% entry in main.tex points at something visible.
	\begin{center}
		{\Large\bfseries Curriculum Vitae}
	\end{center}
	\vspace{0.5em}

	% ------------------------------------------------------------------------
	% Neutralize hand-placed page breaks while reading the resume sections.
	%
	% work_experience.tex contains a bare \newpage that was tuned to the
	% resume's own geometry (0.75in sides, 0.5in top). At the BU margins set
	% above the text block is narrower and taller, so that break no longer
	% falls at the foot of a page -- it fired with about 7 inches of page left,
	% leaving a large blank gap in the middle of the CV.
	%
	% Letting the text flow naturally is what we want here, so \newpage,
	% \clearpage and \pagebreak are turned into harmless \par during the
	% \input calls and restored immediately afterwards. They must be restored,
	% because \end{document} issues a \clearpage of its own.
	%
	% The resume project itself is left completely untouched, so its standalone
	% PDF keeps the pagination you tuned by hand.
	% ------------------------------------------------------------------------
	\let\BUrealNewpage\newpage
	\let\BUrealClearpage\clearpage
	\let\BUrealPagebreak\pagebreak
	\renewcommand{\newpage}{\par}
	\renewcommand{\clearpage}{\par}
	\renewcommand{\pagebreak}{\par}

	% ------------------------------------------------------------------------
	% Section order mirrors main_resume.tex from the resume project.
	% If you add or reorder sections there, mirror the change here.
	% ------------------------------------------------------------------------
	\input{sections/header.tex}
	\input{sections/current_research.tex}
	\input{sections/education.tex}
	\input{sections/work_experience.tex}
	\input{sections/publications.tex}
	\input{sections/poster_presentations.tex}
	\input{sections/lab_skills.tex}
	\input{sections/tech_skills.tex}
	\input{sections/awards.tex}
	\input{sections/languages.tex}
	\input{sections/leadership.tex}

	% ------------------------------------------------------------------------
	% References are commented out ON PURPOSE.
	%
	% Two reasons:
	%  1. Length. With this section the CV runs to five pages, and it is the
	%     only thing on the fifth -- leaving roughly 7 inches of that page
	%     blank. BU guide 1.9 asks candidates to "try to limit the vita to
	%     three or four pages"; without it the CV is exactly four.
	%  2. Convention. Guide 1.9 lists what a vita should contain -- full name,
	%     contact address, prior education, degrees, awards and honours,
	%     professional positions, publications. Referees are not among them,
	%     and a dissertation vita does not normally carry them, since the
	%     document is published openly through ProQuest and OpenBU along with
	%     your referees' phone numbers.
	%
	% Your standalone resume is unaffected and still includes them.
	% To put them back in the dissertation, uncomment the line below and
	% re-run build-cv.sh; expect a five-page CV.
	% ------------------------------------------------------------------------
	% \input{sections/references.tex}

	\let\newpage\BUrealNewpage
	\let\clearpage\BUrealClearpage
	\let\pagebreak\BUrealPagebreak

\end{document}
